\chapter{Thesis Plan}

\section{Working Title}
\label{sec:plan-title}

\textbf{Graph Neural Networks for Decentralized Multi-Agent Reinforcement
Learning in Power-Grid Topology Control}

\section{Central Thesis Question}
\label{sec:plan-question}

\textbf{How can graph neural networks help decentralized
reinforcement learning agents learn reliable power-grid topology control
policies that keep survival performance while scaling to larger environemnt?}
\\
\\
Three terms define what the thesis has to measure:

\begin{itemize}
  \item \textbf{Reliable} is episodic survival on the chronics.
  \item \textbf{Decentralized} means each agent owns a fixed group of
  substations, observes a local subgraph, and acts on its own discrete topology
  actions without seeing what the other agents do. Only the critic is
  centralized.
  \item \textbf{Scaling} means larger grids, with more substations, loads and generators.
\end{itemize}

\section{Research Questions}
\label{sec:plan-research-questions}

\begin{enumerate}
  \item \textbf{How should the grid be represented for a local
  actor?} Is a graph representation useful for learning? If yes, how fine the nodes are, which relations to use, what to add at
  substation level, and how the graph is read out are all design choices.
  Chapter~\ref{ch:methods} defines them and Chapter~\ref{ch:experiments}
  measures which ones change survival.

  % \item \textbf{Can the agents learn to intervene sparsely without losing survival?}
  % Evaluation-time rho overrides, an intervention-gated actor, fixed non-idle
  % penalties, and adaptive Lagrangian intervention budgets are all measured on
  % two axes at once: survival and intervention rate. (Chapter~\ref{ch:sparse-control}.)

  \item \textbf{Is there a transfer possible between a smaller and a larger grid?} Moving from
  \texttt{bus14} to the 36-substation WCCI grid without chainging the weights of a model tests whether a graph encoder
  transfers in practice and not only in parameter shape. 
  (Chapters~\ref{ch:wcci-transfer} and~\ref{ch:screen-e-transfer}.)
\end{enumerate}

\section{Expected Contributions}
\label{sec:plan-contributions}

\begin{itemize}
  \item \textbf{A tuned decentralized MAPPO baseline on \texttt{bus14}.} The
  original framework trained unstably. The training and initialization changes
  that fix it are tested one at a time.

  \item \textbf{An ablation of sparse intervention mechanisms,} reported
  together on survival and intervention frequency. This links the learned
  behavior to what a control room would accept: act rarely, act locally, and
  act when the grid is unsafe.

  \item \textbf{A shared GNN actor and a screened local graph design.} The GNN
  first matches the tuned MLP. Three representations are then compared. Encoder depth and width, message-passing operator, readout, and
  structural augmentations are screened.

  \item \textbf{A model for action scoring that does not
  depend on grid size or action space size.} A candidate-action head scores each action from the
  encoded nodes it changes.
\end{itemize}
